\section{آلية التحكم النبضي (\textenglish{Pulse Train Control})}
\label{sec:pulse_control}

تُعد طريقة التحكم عبر مخرج قطار النبضات (\textenglish{Pulse Train Output - PTO}) الطريقة القياسية والأكثر موثوقية لربط أجهزة الـ \textenglish{PLC} المدمجة مع مشغلات السيرفو في تطبيقات الحركة المستقلة. تعتمد هذه الطريقة على إرسال إشارات رقمية سريعة جداً من المتحكم إلى وحدة السيرفو، حيث يتم فك تشفيرها لفرض كل من الموقع والسرعة.

\subsection{المبدأ الهندسي للتحكم النبضي}
يعتمد جهاز الـ \textenglish{PLC} في تحكمه على عاملين أساسيين عند توليد الإشارة:
\begin{itemize}
    \item \textbf{عدد النبضات الإجمالي (\textenglish{Total Pulse Count}):} يمثل الموقع أو المسافة المقطوعة. كل نبضة يتلقاها السيرفو تُترجم إلى إزاحة ميكانيكية دقيقة (مثلاً، نبضة واحدة = $0.001\text{ mm}$). للوصول إلى هدف محدد، يقوم المتحكم بإصدار عدد دقيق من النبضات ثم يتوقف.
    \item \textbf{تردد النبضات (\textenglish{Pulse Frequency}):} يمثل سرعة المحرك. كلما تقاربت النبضات زمنياً (تردد أعلى)، زادت سرعة دوران المحرك لتنفيذ الأمر. للتحكم في تسارع وتباطؤ الآلة، يقوم الـ \textenglish{PLC} بزيادة أو إنقاص تردد النبضات تدريجياً (\textenglish{Ramp Up / Ramp Down}).
\end{itemize}

\subsection{بروتوكول النبضة والاتجاه (\textenglish{Pulse + Direction})}
من بين عدة بروتوكولات متاحة، يُعد بروتوكول (النبضة + الاتجاه) هو الأسلوب المعتمد في هذا المشروع، نظراً لكفاءته واقتصاده في استخدام مخارج المتحكم. يعتمد هذا البروتوكول على مخرجين رقميين فقط:
\begin{enumerate}
    \item \textbf{مخرج النبضات (\textenglish{Pulse Output}):} يُرسل الإشارة المربعة السريعة جداً.
    \item \textbf{مخرج الاتجاه (\textenglish{Direction Output}):} يُرسل إشارة منطقية ثابتة (إما $0\text{ V}$ أو $24\text{ V}$). إذا كانت الإشارة بوضعية الـ \textenglish{High} يدور المحرك للأمام، وإذا كانت \textenglish{Low} يعكس اتجاهه.
\end{enumerate}



\subsection{المتطلبات العتادية: لماذا مخارج الترانزستور؟}
لتحقيق استجابة ديناميكية سريعة (كما في تطبيق المقص الطائر الذي سيتم استعراضه لاحقاً)، يجب أن تصل سرعة إرسال النبضات إلى مئات الآلاف في الثانية الواحدة. 

لهذا السبب الهندسي، يُمنع تماماً استخدام أجهزة \textenglish{PLC} ذات المخارج الميكانيكية (الريليه \textenglish{Relay}) لأنها بطيئة جداً وتتلف بسرعة. بل يتم الاعتماد حصراً على أجهزة ذات مخارج \textbf{الترانزستور (\textenglish{Transistor Type})}، مثل سلسلة \textenglish{DVP-SV2} من شركة \textenglish{Delta}، والتي تستطيع توليد نبضات بترددات تصل إلى $200\text{ kHz}$، مما يوفر دقة متناهية ونعومة فائقة في حركة محرك السيرفو وتزامنه.